% © 2026 Ing. David Jaroš — CC BY-NC-ND 4.0
%
% This work is licensed under a Creative Commons Attribution-NonCommercial-NoDerivatives
% 4.0 International License (CC BY-NC-ND 4.0).
%
% B_base — Kac-Moody Level k from UBT Action S[Θ] on ψ-Circle
% Track: CORE — α derivation — v67 baseline
% Date: 2026-03-08

\ifdefined\INCLUDEMODE
  % Being included in another document — skip preamble
\else
  \documentclass[11pt,a4paper]{article}
  \usepackage{amsmath,amssymb,amsthm}
  \usepackage{hyperref}
  \usepackage{booktabs}

  \newtheorem{theorem}{Theorem}
  \newtheorem{lemma}{Lemma}
  \newtheorem{proposition}{Proposition}
  \newtheorem{conjecture}{Conjecture}
  \theoremstyle{definition}
  \newtheorem{gap}{Gap}
  \theoremstyle{remark}
  \newtheorem{remark}{Remark}
  \newtheorem{observation}{Observation}

  \title{$B_{\mathrm{base}}$ --- Kac-Moody Level $k$ from UBT Action $\mathcal{S}[\Theta]$\\
    \large Approaches H1 (WZW Normalization), H2 (Absent CS Term),\\
    H3 (Dynkin Index): v67 Baseline}
  \author{Ing.\ David Jaroš}
  \date{2026-03-08 (v67)}

  \begin{document}
  \maketitle
\fi

%--------------------------------------------------------------------
\section{Context and Problem Statement}
\label{sec:context}
%--------------------------------------------------------------------

\noindent\textbf{No circular reasoning.}
$\alpha^{-1} = 137$ and $n^* = 137$ are \emph{never} used as inputs
anywhere in this document.
The only inputs are:
$N_{\mathrm{eff}} = 12$ (proved from $\mathbb{C}\otimes\mathbb{H}$ algebra,
\texttt{step1\_mode\_decomposition.tex} Theorem~1.4),
$\dim_\mathbb{R}(\mathrm{Im}\,\mathbb{H}) = 3$ (algebraic fact), and the UBT
kinetic action $\mathcal{S}[\Theta]$ without explicit coupling constant.

\medskip
\noindent\textbf{Motivation.}
Approach G3 in \texttt{b\_base\_g\_approaches.tex} (\S\ref{sec:context}
above, and Gap~(G3-k) therein) showed that if the $\psi$-circle carries a
$\widehat{\mathrm{SU}(2)}_k$ Kac-Moody symmetry at level $k = 1$, then the
central charge is $c = 1$ and $B_{\mathrm{base}} = c \cdot N_{\mathrm{eff}}^{3/2}
= 41.57$ exactly.  The level $k = 1$ was labelled \textbf{[PARTIAL]}
because no UBT calculation had yet derived $k$ from $\mathcal{S}[\Theta]$.

\medskip
\noindent\textbf{Goal of this document.}
Compute $k$ directly from $\mathcal{S}[\Theta]$ using three independent
approaches:
\begin{itemize}
  \item \textbf{H1:} WZW normalization matching --- read $k$ off the
    coefficient of the kinetic term for $g \in \mathrm{SU}(2)$ on the
    $\psi$-circle.
  \item \textbf{H2:} Absence of Chern-Simons term --- use parity symmetry
    to rule out a CS contribution to $k$ and identify the minimal
    consistent level.
  \item \textbf{H3:} Dynkin index counting --- express $k$ as a sum of
    representation-theoretic indices over the physical modes.
\end{itemize}
If any approach yields $k = 1$ from first principles, Gap~(G3-k)
closes and $B_{\mathrm{base}} = N_{\mathrm{eff}}^{3/2}$ is promoted
from \textbf{[MOTIVATED CONJECTURE]} to \textbf{[PROVED]}.

%--------------------------------------------------------------------
\section{Setup: UBT Action on the $\psi$-Circle}
\label{sec:setup}
%--------------------------------------------------------------------

The UBT kinetic action is (from \texttt{step2\_vacuum\_polarization.tex},
Eq.~(1.1)):
\begin{equation}
  \mathcal{S}_{\mathrm{kin}}[\Theta]
  \;=\;
  \int d^4x\,d\psi\;
  \mathrm{Sc}\!\bigl[(D_\mu\Theta)^\dagger(D^\mu\Theta)\bigr],
  \label{eq:S_kin}
\end{equation}
where $\mathrm{Sc}[\cdot]$ denotes the scalar (identity) part of a
biquaternion, $D_\mu = \partial_\mu + ieA_\mu$ is the electromagnetic
covariant derivative, and $\psi \in [0,2\pi)$ is the compact imaginary-time
circle.  The field $\Theta \in \mathbb{C}\otimes\mathbb{H}$ decomposes as:
\begin{equation}
  \Theta
  \;=\;
  \Theta_0\mathbf{1}
  + \Theta_1\mathbf{i}
  + \Theta_2\mathbf{j}
  + \Theta_3\mathbf{k},
  \qquad \Theta_k \in \mathbb{C},
  \label{eq:decomp}
\end{equation}
with
$\mathrm{Sc}[\Theta^\dagger\Theta]
= |\Theta_0|^2 + |\Theta_1|^2 + |\Theta_2|^2 + |\Theta_3|^2$.
The three charged imaginary components $\Phi_k := \Theta_k$ ($k=1,2,3$)
couple to the electromagnetic field; $\Theta_0$ is neutral and decouples.
The effective action is:
\begin{equation}
  \mathcal{S}_{\mathrm{kin}}^{\mathrm{eff}}
  \;=\;
  \int d^4x\,d\psi\;
  \sum_{k=1}^{3} |D_\mu\Phi_k|^2.
  \label{eq:S_eff}
\end{equation}

The target value of $B_{\mathrm{base}}$ (from $N_{\mathrm{eff}} = 12$,
proved independently of $\alpha^{-1}$) is:
\begin{equation}
  B_{\mathrm{base}}
  \;=\;
  N_{\mathrm{eff}}^{3/2}
  \;=\;
  12^{3/2}
  \;=\;
  24\sqrt{3}
  \;\approx\;
  41.569.
  \label{eq:B_base_target}
\end{equation}
All numerical values $\approx 41.57$ below refer to this exact algebraic expression.

\subsection{SU(2) identification on the $\psi$-circle}
\label{sec:SU2_id}

The Lie algebra $\mathfrak{su}(2)$ is spanned by $\{\mathbf{i},\mathbf{j},\mathbf{k}\}$
with $[\mathbf{i},\mathbf{j}] = 2\mathbf{k}$ etc., and the compact group
$\mathrm{SU}(2)$ corresponds to the \emph{real} unit quaternions $S^3$:
\begin{equation}
  \mathrm{SU}(2)
  \;\cong\;
  \{g = a_0\mathbf{1} + a_1\mathbf{i} + a_2\mathbf{j} + a_3\mathbf{k}
   \;|\; a_\mu \in \mathbb{R},\; a_0^2 + a_1^2 + a_2^2 + a_3^2 = 1\}
  \;\subset\; \mathbb{H}
  \;\subset\; \mathbb{C}\otimes\mathbb{H}.
  \label{eq:SU2_def}
\end{equation}
For the WZW identification, we restrict to real-valued field configurations
$\Theta_{\mathrm{real}}(x,\psi) \in \mathbb{H}$ and decompose:
\begin{equation}
  \Theta_{\mathrm{real}} \;=\; r \cdot g,
  \qquad
  r = |\Theta_{\mathrm{real}}| \in \mathbb{R}_{\geq 0},
  \quad
  g \in \mathrm{SU}(2).
  \label{eq:polar}
\end{equation}
This decomposition identifies the ``radial'' mode $r$ and the ``angular'' mode
$g$ on the $\psi$-circle.  The biquaternionic field $\Theta \in \mathbb{C}\otimes\mathbb{H}$
does \emph{not} generically lie in $\mathrm{SU}(2)$; the restriction to the
real locus is a necessary assumption for the WZW identification (see
Remark~\ref{rem:complex_obstruction}).

\begin{remark}[Obstruction from complex values]
  \label{rem:complex_obstruction}
  The full UBT field $\Theta$ has $\Theta_k \in \mathbb{C}$,
  so $\Theta \in \mathbb{C}\otimes\mathbb{H}$ with $8$ real degrees of freedom.
  The real locus $\mathbb{H} \subset \mathbb{C}\otimes\mathbb{H}$ has only $4$
  real degrees of freedom; the projection $\Theta \to \mathrm{Re}(\Theta) \in \mathbb{H}$
  discards the imaginary parts $\mathrm{Im}(\Theta_k)$.
  Unit-norm elements of $\mathbb{C}\otimes\mathbb{H}$ form $\mathrm{SL}(2,\mathbb{C})$,
  not $\mathrm{SU}(2)$.
  Approaches H1--H3 therefore apply only under the restriction to the
  $\mathrm{SU}(2)$ locus, and this restriction is not derived from UBT axioms.
\end{remark}

%--------------------------------------------------------------------
\section{Approach H1: WZW Normalization}
\label{sec:H1}
%--------------------------------------------------------------------

\textbf{Status: [DEAD END] --- the vacuum expectation value
$\langle|\Theta|^2\rangle_{\mathrm{vac}}$ needed to extract $k$ either depends
on $n^* = 137$ (circular) or reduces to a normalization convention.
Neither path derives $k = 1$ from UBT axioms alone.}

\subsection{Deriving $k$ from the kinetic action}
\label{sec:H1_derivation}

Under the polar decomposition~\eqref{eq:polar}, the kinetic action restricted
to $\psi$-circle dependence (suppressing the $d^4x$ integral at fixed
$\langle r^2\rangle = r_{\mathrm{vac}}^2$) becomes:
\begin{equation}
  \mathcal{S}_{\mathrm{kin}}[\psi]
  \;=\;
  r_{\mathrm{vac}}^2 \int_0^{2\pi} d\psi\;
  \mathrm{Sc}\bigl[(g^\dagger \partial_\psi g)^\dagger\,(g^\dagger \partial_\psi g)\bigr].
  \label{eq:S_psi}
\end{equation}

\noindent\textit{Why this form:} Inserting $\Theta = r_{\mathrm{vac}} g$ with
$\partial_\psi r_{\mathrm{vac}} = 0$ (we evaluate at the vacuum), and using
$\partial_\psi g^\dagger = -\omega^\dagger g^\dagger = \omega g^\dagger$
where $\omega = g^\dagger\partial_\psi g \in \mathfrak{su}(2)$
(antihermitian: $\omega^\dagger = -\omega$), one obtains
$(g^\dagger\partial_\psi g)^\dagger = -\omega$ and thus:
\begin{equation}
  \mathrm{Sc}\bigl[(g^\dagger\partial_\psi g)^\dagger\,(g^\dagger\partial_\psi g)\bigr]
  \;=\;
  \mathrm{Sc}[-\omega^2]
  \;=\;
  |\omega|^2
  \;\equiv\;
  \omega_1^2 + \omega_2^2 + \omega_3^2,
  \label{eq:Sc_omega}
\end{equation}
where $\omega = \omega_1\mathbf{i} + \omega_2\mathbf{j} + \omega_3\mathbf{k}$
and we used $\mathbf{e}^2 = -1$ for any $\mathbf{e} \in \{\mathbf{i},\mathbf{j},\mathbf{k}\}$
so that $\omega^2 = -(|\omega|^2)\mathbf{1}$ (no cross terms because
$\{\mathbf{e}_a, \mathbf{e}_b\} = -2\delta_{ab}$ in $\mathbb{H}$).
Hence:
\begin{equation}
  \mathcal{S}_{\mathrm{kin}}[\psi]
  \;=\;
  r_{\mathrm{vac}}^2 \int_0^{2\pi} d\psi\; |\omega|^2.
  \label{eq:S_kinematic_psi}
\end{equation}

\subsection{Standard WZW normalization}
\label{sec:WZW_norm}

The Wess-Zumino-Witten action on $\mathrm{SU}(2)$ at level $k$ on the circle
$S^1$ is (in the kinetic-term form, without the topological Wess-Zumino term
which integrates over a 3-manifold bounding $S^1$):
\begin{equation}
  \mathcal{S}_{\mathrm{WZW}}^{(k)}
  \;=\;
  -\frac{k}{4\pi}
  \int_0^{2\pi} d\psi\;
  \mathrm{Tr}_{\mathrm{fund}}\!\bigl[(g^\dagger\partial_\psi g)^2\bigr],
  \label{eq:WZW}
\end{equation}
where $\mathrm{Tr}_{\mathrm{fund}}$ is the matrix trace in the fundamental
($2$-dimensional) representation of $\mathrm{SU}(2)$.
For $\omega = \omega_a(i\sigma_a/2) \in \mathfrak{su}(2)$
(in terms of Pauli matrices $\sigma_a$):
\begin{equation}
  \mathrm{Tr}_{\mathrm{fund}}(\omega^2)
  \;=\;
  \mathrm{Tr}_{\mathrm{fund}}\!\left[
    -\frac{\omega_1^2+\omega_2^2+\omega_3^2}{4}\,I_2
  \right]
  \;=\;
  -\frac{|\omega|^2}{2}.
  \label{eq:Tr_omega2}
\end{equation}
Substituting into \eqref{eq:WZW}:
\begin{equation}
  \mathcal{S}_{\mathrm{WZW}}^{(k)}
  \;=\;
  \frac{k}{8\pi}
  \int_0^{2\pi} d\psi\; |\omega|^2.
  \label{eq:WZW_positive}
\end{equation}

\subsection{Level identification}
\label{sec:level_id}

Matching \eqref{eq:S_kinematic_psi} against \eqref{eq:WZW_positive}
coefficient-by-coefficient:
\begin{equation}
  \boxed{
    k
    \;=\;
    8\pi\, r_{\mathrm{vac}}^2
    \;=\;
    8\pi\, \langle|\Theta_{\mathrm{real}}|^2\rangle_{\mathrm{vac}}.
  }
  \label{eq:k_H1}
\end{equation}
\textit{Why the vacuum amplitude matters:} The WZW level $k$ must be a
positive integer (quantisation condition from the topological Wess-Zumino
term); Eq.~\eqref{eq:k_H1} says $k$ equals $8\pi$ times the
squared amplitude of the $\psi$-circle configuration at the vacuum.

\subsection{Can $r_{\mathrm{vac}}^2$ be determined from UBT without $n^*$?}
\label{sec:r_vac}

There are two natural candidates for $r_{\mathrm{vac}}^2$:

\paragraph{Candidate 1 — from the effective potential minimum.}
The UBT effective potential $V_{\mathrm{eff}}(n)$ has its minimum at $n^*$.
In principle, the vacuum amplitude $r_{\mathrm{vac}}^2 = \langle|\Theta|^2\rangle$
could be extracted from the equations of motion at this minimum.
However, $n^*$ is the \emph{winding number} (phase variable), not the
radial amplitude $r$; the two are \emph{conjugate} variables on the
$\psi$-circle (angle vs.\ radius).
More concretely, the minimum $n^* = 137$ is the output we are trying
to derive ($n^* \leftrightarrow \alpha^{-1}$), and any extraction of
$r_{\mathrm{vac}}$ from $V_{\mathrm{eff}}$ at $n = n^*$ would use $n^*$
as input.

\begin{observation}[H1 Circularity]
  \label{obs:H1_circular}
  If $r_{\mathrm{vac}}^2$ is computed from the vacuum structure of
  $V_{\mathrm{eff}}(n)$ at $n = n^*$, the result depends on $n^*$,
  which equals $\alpha^{-1} = 137$.
  This makes Eq.~\eqref{eq:k_H1} circular:
  the very quantity ($\alpha^{-1}$) that $k = 1$ is supposed to derive
  (via $B_{\mathrm{base}} = c_{k=1}\cdot N_{\mathrm{eff}}^{3/2}$) would
  already appear as an input to $r_{\mathrm{vac}}$.
\end{observation}

\paragraph{Candidate 2 — canonical normalization convention.}
The UBT action~\eqref{eq:S_kin} has no explicit coupling-constant prefactor.
A natural choice is to impose the canonical normalization
$\mathrm{Sc}[\Theta^\dagger\Theta] = 1$ (dimensionless, in natural units
$\hbar = c = 1$), giving $r_{\mathrm{vac}}^2 = 1$.
Then from Eq.~\eqref{eq:k_H1}:
\begin{equation}
  k
  \;=\;
  8\pi \times 1
  \;=\;
  8\pi
  \;\approx\;
  25.1,
  \label{eq:k_convention}
\end{equation}
which is not a positive integer.
Alternatively, rescaling the action as
$\tilde{\mathcal{S}}_{\mathrm{kin}} = (1/4\pi)\mathcal{S}_{\mathrm{kin}}$
and $r_{\mathrm{vac}}^2 = 1/(4\pi)$ gives $k = 2$.
No canonical procedure produces $k = 1$ from the normalization convention alone.

\begin{observation}[H1 normalization non-uniqueness]
  \label{obs:H1_norm}
  Under the convention $r_{\mathrm{vac}}^2 = 1$,
  Eq.~\eqref{eq:k_H1} gives $k = 8\pi \notin \mathbb{Z}$.
  The level $k = 1$ requires $r_{\mathrm{vac}}^2 = 1/(8\pi)$, which
  is not distinguished by any algebraic property of
  $\mathbb{C}\otimes\mathbb{H}$.
\end{observation}

\textbf{H1 result summary:}
\begin{itemize}
  \item $k = 8\pi\, r_{\mathrm{vac}}^2$ (exact formula from kinetic action matching).
  \item $r_{\mathrm{vac}}^2$ from $V_{\mathrm{eff}}$ minimum: circular ($n^*$ is an output,
    not an input).
  \item $r_{\mathrm{vac}}^2 = 1$ (canonical normalization): $k = 8\pi \notin \mathbb{Z}$.
  \item $k = 1$: requires $r_{\mathrm{vac}}^2 = 1/(8\pi)$, not algebraically
    natural.
  \item \textbf{[DEAD END]:} H1 does not derive $k$ from UBT axioms without
    a free normalization parameter or circular input.
\end{itemize}

%--------------------------------------------------------------------
\section{Approach H2: Absence of Chern-Simons Term}
\label{sec:H2}
%--------------------------------------------------------------------

\textbf{Status: [PARTIAL] --- parity symmetry rigorously excludes
a Chern-Simons contribution to $k$.
Combined with the free-boson realisation of the WZW model, this
identifies $k = 1$ as the \emph{unique minimal consistent level} for
UBT.
However, ``minimal'' is a physical argument, not a theorem from
UBT axioms; $k > 1$ from a non-CS mechanism cannot be excluded by H2 alone.}

\subsection{Structure of the UBT action}
\label{sec:action_structure}

The full UBT action is:
\begin{equation}
  \mathcal{S}[\Theta]
  \;=\;
  \mathcal{S}_{\mathrm{kin}}[\Theta]
  + \mathcal{S}_{\mathrm{pot}}[\Theta],
  \label{eq:S_full}
\end{equation}
where $\mathcal{S}_{\mathrm{kin}}$ is the kinetic term~\eqref{eq:S_kin}
and $\mathcal{S}_{\mathrm{pot}}$ is the effective potential derived from
one-loop vacuum polarisation.
Both terms are built from $\mathrm{Sc}[\Theta^\dagger\Theta]$ and its
derivatives; no topological density of Chern-Simons type
$\epsilon^{\mu\nu\rho}\mathrm{Tr}(A_\mu\partial_\nu A_\rho + \tfrac{2}{3}A_\mu A_\nu A_\rho)$
appears in either term.

\subsection{Parity argument for the absence of CS}
\label{sec:parity_CS}

\begin{proposition}[No Chern-Simons term in $\mathcal{S}[\Theta]$]
  \label{prop:no_CS}
  The action $\mathcal{S}[\Theta]$ contains no Chern-Simons term.
\end{proposition}

\begin{proof}
  A Chern-Simons term $\mathcal{S}_{\mathrm{CS}} \propto
  \int \epsilon^{\mu\nu\rho}\mathrm{Tr}(A_\mu\partial_\nu A_\rho + \cdots)$
  is odd under the parity transformation $P$: it changes sign when
  spatial coordinates are reflected ($x^i \to -x^i$), because the
  Levi-Civita symbol $\epsilon^{\mu\nu\rho}$ is a pseudotensor.

  UBT has an explicit $P$-even/odd decomposition of the $\psi$-circle
  into the primary sector ($n^* = 137$) and the mirror sector
  ($n^* = 139$, $CP$-conjugate), where $P$ is a global symmetry of the
  two-sector system (see \texttt{step1\_psi\_parity.tex}).
  Since the full theory is $P$-symmetric, no $P$-odd term can appear in
  $\mathcal{S}[\Theta]$.

  Furthermore, the explicit construction of $\mathcal{S}[\Theta]$ from
  the biquaternionic norm $\mathrm{Sc}[\Theta^\dagger\Theta]$ shows that
  every term is $P$-even.
  \textit{(Why: $\mathrm{Sc}[\Theta^\dagger\Theta] = \sum_k|\Theta_k|^2$
  is a scalar density under all discrete symmetries; its Lagrangian
  density transforms as a genuine scalar under $P$.)}
\end{proof}

\begin{remark}
  Proposition~\ref{prop:no_CS} is a \emph{structural} result: it follows
  from the algebra, not from the specific value of any parameter.
  It holds for all values of $n^*$, in particular without requiring
  $n^* = 137$.
\end{remark}

\subsection{The free-boson realisation and the minimal level}
\label{sec:free_boson}

The $\widehat{\mathrm{SU}(2)}_k$ WZW model at level $k = 1$ is
\emph{equivalent} to a free Dirac fermion on the circle (via bosonisation)
or, equivalently, to three free massless real scalars compactified at the
$\mathrm{SU}(2)_1$ self-dual radius.
This is the standard textbook result sometimes called the
``free-fermion point'' or ``Gepner point'' of the WZW model.

More generally, for a non-linear sigma model on $\mathrm{SU}(2)$ with no
Chern-Simons term:
\begin{itemize}
  \item The WZW term (Wess-Zumino topological term) is generated at one
    loop with level $k = 1$ per fundamental scalar degree of freedom.
  \item Without a CS term, the only source of a positive integer shift to
    $k$ is the number of fundamental scalars contributing to the one-loop
    WZW anomaly.
  \item A single free real scalar on $\mathrm{SU}(2)$ generates $k = 1$.
  \item The minimal consistent level for any unitary $\widehat{\mathrm{SU}(2)}$
    WZW model is $k = 1$ (positivity of the central charge requires
    $k \geq 1$).
\end{itemize}

\begin{observation}[H2: Minimal level for UBT]
  \label{obs:H2_minimal}
  Since Proposition~\ref{prop:no_CS} rules out a CS term, the UBT
  $\psi$-circle WZW level $k$ is determined entirely by the kinetic action.
  The kinetic action~\eqref{eq:S_kin} describes a \emph{free minimally-coupled
  scalar} $\Theta$ (no higher derivatives, no self-interaction in the kinetic
  term).
  By the free-boson/WZW correspondence, this corresponds to the minimal
  level $k = 1$.
\end{observation}

\subsection{Limitation of H2}
\label{sec:H2_limit}

The free-boson/WZW correspondence in Observation~\ref{obs:H2_minimal}
strictly applies when $\Theta$ is a \emph{real} scalar on $\mathrm{SU}(2)$.
As noted in Remark~\ref{rem:complex_obstruction}, the UBT field has
$\Theta_k \in \mathbb{C}$; the correct group manifold for complex biquaternions
of unit norm is $\mathrm{SL}(2,\mathbb{C})$, not $\mathrm{SU}(2)$.
The WZW model on $\mathrm{SL}(2,\mathbb{C})$ is a non-compact WZNW model
(Wess-Zumino-Novikov-Witten) whose level structure differs from the
$\mathrm{SU}(2)$ case.

Therefore:
\begin{itemize}
  \item H2 \emph{proves} that no CS term shifts $k$ upward.
  \item H2 \emph{motivates} $k = 1$ as the minimal level by the
    free-boson correspondence, but only rigorously for the real locus
    $\Theta \in \mathbb{H} \subset \mathbb{C}\otimes\mathbb{H}$.
  \item H2 does \emph{not} prove that $k = 1$ rather than $k > 1$, because a
    higher level could arise from the kinetic term normalisation (H1) without
    any CS contribution, or from the non-compact ($\mathrm{SL}(2,\mathbb{C})$)
    nature of the full field.
\end{itemize}

\textbf{H2 result summary:}
\begin{itemize}
  \item CS term absent: proved from $P$-symmetry.
  \item Minimal consistent level without CS: $k = 1$ (free-boson correspondence).
  \item Restriction: free-boson correspondence is rigorous only on the real
    $\mathrm{SU}(2)$ locus; the full complex-valued UBT field could have
    $k > 1$ from the kinetic normalisation.
  \item \textbf{[PARTIAL]:} H2 proves absence of CS and identifies $k = 1$
    as the natural minimal level, but does not uniquely force $k = 1$ for
    the full $\mathbb{C}\otimes\mathbb{H}$ theory.
\end{itemize}

%--------------------------------------------------------------------
\section{Approach H3: Dynkin Index Counting}
\label{sec:H3}
%--------------------------------------------------------------------

\textbf{Status: [DEAD END] --- the natural representation-theoretic
counting of the $N_{\mathrm{eff}} = 12$ modes gives $k = 3/2$
(half-integer), which yields
$c_{\widehat{\mathrm{SU}(2)},3/2} = 9/7 \approx 1.286$ and
$B_{\mathrm{base}} \approx 53.5 \neq 41.57$.
Alternative representations (adjoint, $k=1$ from two modes) also fail.}

\subsection{Dynkin index formula for the WZW level}
\label{sec:H3_formula}

In the operator-algebraic derivation of WZW models, the level $k$ of
$\widehat{\mathfrak{g}}_k$ equals the sum of Dynkin indices over all
field modes contributing to the current algebra:
\begin{equation}
  k
  \;=\;
  \sum_i n_i \cdot T(\mathrm{rep}_i),
  \label{eq:k_Dynkin}
\end{equation}
where $n_i$ is the multiplicity of representation $\mathrm{rep}_i$ and
$T(\mathrm{rep})$ is the Dynkin index (normalised so that the fundamental
representation of $\mathrm{SU}(N)$ has $T = 1/2$).

For $\mathrm{SU}(2)$: $T(\mathbf{1}) = 0$, $T(\mathbf{2}) = 1/2$,
$T(\mathbf{3}) = 2$, $T(\mathbf{j+1}) = j(j+1)(2j+1)/3$.

\subsection{The three imaginary components as $\mathrm{SU}(2)$ representations}
\label{sec:H3_reps}

The three imaginary quaternion directions
$\{\mathbf{i},\mathbf{j},\mathbf{k}\} \cong \mathrm{Im}(\mathbb{H}) \cong \mathbb{R}^3$
form the adjoint representation $\mathbf{3}$ of $\mathrm{SU}(2)$.
However, each $\Phi_k \in \mathbb{C}$ carries a \emph{complex} degree of
freedom.
A complex scalar in the fundamental $\mathbf{2}$ has $T = 1/2$; a real
scalar in the $\mathbf{3}$ has $T = 2$; a complex scalar in the $\mathbf{3}$
has $T = 2$.

\paragraph{Case A --- Three fundamental modes ($N_{\mathrm{phases}} = 3$, each in $\mathbf{2}$):}
\begin{equation}
  k_A
  \;=\;
  3 \times T(\mathbf{2})
  \;=\;
  3 \times \tfrac{1}{2}
  \;=\;
  \tfrac{3}{2}.
  \label{eq:k_A}
\end{equation}
For $k = 3/2$ (half-integer levels are allowed in $\widehat{\mathrm{SU}(2)}$):
\begin{equation}
  c_{\widehat{\mathrm{SU}(2)},\,3/2}
  \;=\;
  \frac{3 \times 3/2}{3/2 + 2}
  \;=\;
  \frac{9/2}{7/2}
  \;=\;
  \frac{9}{7}
  \;\approx\;
  1.286.
  \label{eq:c_k32}
\end{equation}
\begin{equation}
  B_{\mathrm{base}}(k=3/2)
  \;=\;
  c_{\widehat{\mathrm{SU}(2)},\,3/2} \times N_{\mathrm{eff}}^{3/2}
  \;=\;
  \tfrac{9}{7} \times 41.57
  \;\approx\;
  53.4.
  \label{eq:B_k32}
\end{equation}
This is $28\%$ above the target $41.57$.
\textbf{[DEAD END] for $k_A = 3/2$.}

\paragraph{Case B --- One adjoint mode ($\mathbf{3}$):}
\begin{equation}
  k_B
  \;=\;
  1 \times T(\mathbf{3})
  \;=\;
  2.
  \label{eq:k_B}
\end{equation}
\begin{equation}
  c_{\widehat{\mathrm{SU}(2)},\,2}
  \;=\;
  \frac{3 \times 2}{2 + 2}
  \;=\;
  \frac{6}{4}
  \;=\;
  \frac{3}{2}.
  \label{eq:c_k2}
\end{equation}
\begin{equation}
  B_{\mathrm{base}}(k=2)
  \;=\;
  \tfrac{3}{2} \times 41.57
  \;\approx\;
  62.4.
  \label{eq:B_k2}
\end{equation}
\textbf{[DEAD END] for $k_B = 2$.}

\paragraph{Case C --- $N_{\mathrm{eff}} = 12$ modes, each in $\mathbf{2}$:}
\begin{equation}
  k_C
  \;=\;
  12 \times T(\mathbf{2})
  \;=\;
  12 \times \tfrac{1}{2}
  \;=\;
  6.
  \label{eq:k_C}
\end{equation}
\begin{equation}
  c_{\widehat{\mathrm{SU}(2)},\,6}
  \;=\;
  \frac{3 \times 6}{6 + 2}
  \;=\;
  \frac{18}{8}
  \;=\;
  \frac{9}{4}
  \;=\;
  2.25.
  \label{eq:c_k6}
\end{equation}
\begin{equation}
  B_{\mathrm{base}}(k=6)
  \;=\;
  2.25 \times 41.57
  \;\approx\;
  93.5.
  \label{eq:B_k6}
\end{equation}
\textbf{[DEAD END] for $k_C = 6$.}

\paragraph{Could $k = 1$ arise from H3?}
For $k = 1$ from Eq.~\eqref{eq:k_Dynkin}, we need:
$\sum_i n_i T(\mathrm{rep}_i) = 1$.
The natural options:
\begin{itemize}
  \item $n = 2$, $T = 1/2$ (two fundamental modes): $k = 1$.
    But the UBT derivation fixes $N_{\mathrm{phases}} = 3$ (proved algebraically
    from $\dim_\mathbb{R}(\mathrm{Im}\,\mathbb{H}) = 3$), not $2$.
    Reducing to $2$ modes would contradict the proved value.
  \item $n = 1$, $T = 1$: requires a spin-$j$ representation with
    $T(j) = j(j+1)(2j+1)/3 = 1$, giving $j \approx 0.77$ --- not a
    half-integer. No such representation exists in $\mathrm{SU}(2)$.
\end{itemize}
No natural combination of $\mathrm{SU}(2)$ representations with
$N_{\mathrm{phases}} = 3$ modes gives $k = 1$ exactly.

\begin{center}
\begin{tabular}{lllll}
  \toprule
  Representation & $n$ modes & $T(\mathrm{rep})$ & $k$ & $B_{\mathrm{base}}$ \\
  \midrule
  $\mathbf{2}$ (fundamental) & 3 & $1/2$ & $3/2$ & $53.4$ \quad [DEAD END] \\
  $\mathbf{3}$ (adjoint) & 1 & $2$ & $2$ & $62.4$ \quad [DEAD END] \\
  $\mathbf{2}$ (fundamental) & 12 & $1/2$ & $6$ & $93.5$ \quad [DEAD END] \\
  $\mathbf{2}$ (fundamental) & 2 & $1/2$ & $1$ & $41.57$ \quad [CONTRADICTS $N_{\mathrm{phases}}=3$] \\
  \bottomrule
\end{tabular}
\end{center}

\textbf{H3 result summary:}
\begin{itemize}
  \item Most natural counting ($N_{\mathrm{phases}} = 3$, fund.\ rep.):
    $k = 3/2$, $B_{\mathrm{base}} \approx 53.4 \neq 41.57$.
  \item Adjoint rep: $k = 2$, $B_{\mathrm{base}} \approx 62.4 \neq 41.57$.
  \item $k = 1$ requires $n = 2$ fundamental modes, contradicting the
    proved $N_{\mathrm{phases}} = 3$.
  \item \textbf{[DEAD END]:} No natural $\mathrm{SU}(2)$-representation
    assignment to the UBT modes with $N_{\mathrm{phases}} = 3$ gives $k = 1$.
\end{itemize}

%--------------------------------------------------------------------
\section{Synthesis and Conclusion}
\label{sec:synthesis}
%--------------------------------------------------------------------

\subsection{Comparison of H1--H3}

\begin{center}
\begin{tabular}{lll}
  \toprule
  Approach & Result & Status \\
  \midrule
  H1: WZW normalization
    & $k = 8\pi r_{\mathrm{vac}}^2$; $r_{\mathrm{vac}}$ unknown
    & \textbf{[DEAD END]} \\
  H1 (variant): canonical $r^2 = 1$
    & $k = 8\pi \notin \mathbb{Z}$
    & \textbf{[DEAD END]} \\
  H2: Absence of CS term
    & No CS shift; $k=1$ is minimal free-scalar level
    & \textbf{[PARTIAL]} \\
  H3: $N_{\mathrm{phases}}=3$ fundamental modes
    & $k = 3/2$; $B_{\mathrm{base}} = 53.4$
    & \textbf{[DEAD END]} \\
  H3: adjoint ($\mathbf{3}$)
    & $k = 2$; $B_{\mathrm{base}} = 62.4$
    & \textbf{[DEAD END]} \\
  H3: 2 fundamental modes
    & $k = 1$; contradicts $N_{\mathrm{phases}}=3$
    & \textbf{[DEAD END]} \\
  \bottomrule
\end{tabular}
\end{center}

\subsection{Is $k = 1$ proved?}

\textbf{Short answer: No.}
None of H1--H3 derives $k = 1$ from UBT first principles without either
(a) introducing a normalization convention (H1), (b) relying on the
real-locus approximation $\Theta \in \mathbb{H}$ (H2), or (c) contradicting
the proved value $N_{\mathrm{phases}} = 3$ (H3).

\textbf{Best positive result:}
H2 (absence of Chern-Simons) establishes that no integer uplift of $k$
comes from a CS term, and the free-boson correspondence identifies
$k = 1$ as the unique minimal level for a free scalar on the $\mathrm{SU}(2)$
locus.
This strengthens the motivation for $k = 1$ beyond the original G3
observation, but does not constitute a proof.

\textbf{Key obstruction:}
The full UBT field is complex-valued ($\Theta_k \in \mathbb{C}$), and unit-norm
complex biquaternions form $\mathrm{SL}(2,\mathbb{C})$ rather than $\mathrm{SU}(2)$.
The WZW theory for $\mathrm{SL}(2,\mathbb{C})$ is a non-unitary, non-compact
model; the standard level-quantisation argument does not apply.
Resolving this obstruction — either by projecting to $\mathrm{SU}(2)$ with
a rigorous justification, or by analysing the non-compact WZW model for
$\mathrm{SL}(2,\mathbb{C})$ --- is the concrete next step.

%--------------------------------------------------------------------
\section{Updated Gap Status}
\label{sec:gaps}
%--------------------------------------------------------------------

\begin{gap}[Gap (G3-k): Level of $\widehat{\mathrm{SU}(2)}$ in UBT — updated]
  \label{gap:G3_k_updated}
  \textbf{Previous status (b\_base\_g\_approaches.tex):} [OPEN]

  \textbf{After H1--H3:} [PARTIAL PROGRESS]

  \begin{itemize}
    \item H2 proves: no Chern-Simons term, hence no CS contribution to $k$.
    \item H2 motivates: $k = 1$ is the minimal consistent level for the
      $\mathrm{SU}(2)$ locus of the UBT field.
    \item H1 shows: $k$ is not uniquely determined from the kinetic
      action without knowing $r_{\mathrm{vac}}$, which is either circular
      (from $n^*$) or a free normalization convention.
    \item H3 shows: the natural Dynkin-index counting with $N_{\mathrm{phases}}=3$
      gives $k = 3/2$ (wrong $B_{\mathrm{base}}$), not $k = 1$.
    \item Key open question: Can the restriction $\Theta \in \mathrm{SU}(2)$
      (real locus) be justified from UBT axioms?
      If yes, H2 promotes $k = 1$ to [MOTIVATED ARGUMENT].
  \end{itemize}
  Status: \textbf{[PARTIAL --- no-CS proved; $k=1$ minimal but not uniquely forced]}.
\end{gap}

\begin{gap}[Gap (H-complex): WZW theory for $\mathrm{SL}(2,\mathbb{C})$ on $\psi$-circle]
  \label{gap:H_complex}
  The correct WZW/WZNW analysis for unit-norm complex biquaternions uses
  the non-compact group $\mathrm{SL}(2,\mathbb{C})$ rather than $\mathrm{SU}(2)$.
  The representation theory of $\widehat{\mathfrak{sl}}(2,\mathbb{C})_k$
  (non-compact WZW model) allows non-integer and even complex levels $k$.
  Determining the level for the UBT $\psi$-circle action in the
  non-compact setting is the next concrete direction.
  Status: \textbf{[OPEN --- new gap identified by H1--H3 analysis]}.
\end{gap}

\medskip
\noindent\textbf{Impact on $B_{\mathrm{base}}$ status:}
The overall status of $B_{\mathrm{base}} = N_{\mathrm{eff}}^{3/2}$ remains
\textbf{[MOTIVATED CONJECTURE]}.
The H2 result (no CS term + minimal level $k = 1$ on the $\mathrm{SU}(2)$
locus) provides an additional motivation alongside Approach A2
(Gaussian path integral) and G3 ($k=1$ matches $B_{\mathrm{base}}$ exactly),
but does not close the derivation.

%--------------------------------------------------------------------
\section{In Plain Language}
\label{sec:plain}
%--------------------------------------------------------------------

\textbf{In plain language:}
We tried three methods to determine the ``level'' $k$ --- a positive
integer that labels how strong a symmetry algebra is on the imaginary-time
circle of UBT.
The reason $k$ matters: if $k = 1$, then the Kac-Moody central charge
equals exactly $1$, and the formula $B_{\mathrm{base}} = c \cdot N_{\mathrm{eff}}^{3/2}$
gives the correct value $41.57$ with no free parameters.
If $k \neq 1$, the formula fails.

\textbf{Method H1 (WZW normalization):}
We computed the coefficient in front of the kinetic energy of the group-valued
part of $\Theta$ on the imaginary-time circle, and compared it to the
standard formula for the $k$-th level WZW model.
The result is: $k$ equals $8\pi$ times the squared ``amplitude'' of the
field at the vacuum.
The problem is that we do not know this vacuum amplitude independently:
reading it off the vacuum structure of the effective potential would
require knowing $\alpha^{-1} = 137$ in advance (circular),
and the ``natural'' normalization ($|\Theta|^2 = 1$) gives $k = 8\pi \approx 25$,
which is not an integer.
\textit{Dead end.}

\textbf{Method H2 (No Chern-Simons term):}
We showed --- rigorously --- that the UBT action contains no
Chern-Simons term.
Such a term is ``P-odd'' (it changes sign when you look at it in a
mirror), but UBT has a mirror sector (the twin-prime sector at $n^*=139$),
so the full theory is mirror-symmetric and a CS term is forbidden.
Without a CS term, the level $k$ receives no extra integer contribution.
The known physics of free scalar fields on the sphere $S^3$ tells us
that the minimal consistent level is $k = 1$.
So H2 gives a strong physical argument for $k = 1$ (no extra contributions +
free scalar gives minimal level), but it is not a full proof:
the UBT field is \emph{complex}-valued, and the free-scalar/$k=1$ argument
strictly applies only to \emph{real}-valued fields on $S^3$.
\textit{Partial result: motivates $k = 1$, but does not prove it.}

\textbf{Method H3 (Counting representations):}
We tried to compute $k$ by adding up the ``weight'' (Dynkin index) of
each UBT field mode.
UBT has $N_{\mathrm{phases}} = 3$ imaginary components, each transforming
under $\mathrm{SU}(2)$.
If each transforms as the fundamental (spin-$\tfrac{1}{2}$) representation,
we get $k = 3 \times \tfrac{1}{2} = \tfrac{3}{2}$.
The central charge at $k = \tfrac{3}{2}$ is $\tfrac{9}{7}$, which gives
$B_{\mathrm{base}} \approx 53.4$ --- $28\%$ too large.
The only way to get $k = 1$ from this counting is to use $n = 2$
(not $3$) modes, contradicting the proved value $N_{\mathrm{phases}} = 3$.
\textit{Dead end.}

\textbf{Overall verdict:}
After three new approaches, $k = 1$ remains well-motivated but not proved.
The strongest new result is the rigorous \emph{absence} of a Chern-Simons
term (H2), which removes one potential obstruction and focuses the search.
The key open question is whether the complex-valued UBT field effectively
reduces to a real-valued field on the $\mathrm{SU}(2) = S^3$ submanifold
of $\mathbb{C}\otimes\mathbb{H}$; if yes, $k = 1$ follows from standard
free-scalar/WZW theory.

\ifdefined\INCLUDEMODE\else
\end{document}
\fi
